\subsection{Prenotare un test drive}\label{sec:uc-test-drive}
\begin{xltabular}{\textwidth}{| >{\bfseries}p{.22\textwidth} | >{\raggedright\arraybackslash}X |}
\caption{Caso d'uso: Prenotare un test drive}\label{tab:uc-test-drive}\\
\hline
\rowcolor{gray!15} Campo & \textbf{Contenuto}\\\hline
\endfirsthead
\multicolumn{2}{c}{{\tablename\ \thetable{} -- Continua dalla pagina precedente}}\\
\hline
\rowcolor{gray!15} Campo & \textbf{Contenuto}\\\hline
\endhead
\hline
\multicolumn{2}{r}{\textit{Continua nella pagina successiva...}}\\
\endfoot
\endlastfoot
Nome & Prenotare un test drive\\\hline
Livello & Obiettivo Utente\\\hline
Descrizione & Il Customer prenota una prova di guida di un'ora per un veicolo disponibile nel catalogo.\\\hline
Pre-condizioni & Il Customer ha effettuato l'accesso; il veicolo esiste ed è in stato \code{SALE} o \code{RENT}; la data e la fascia oraria scelte non sono passate.\\\hline
Post-condizioni & Una prenotazione di test drive è registrata per il Customer, il veicolo e lo slot selezionato.\\\hline
Attori & Customer\\\hline
Test & \code{UserGoalsFunctionalTest.bookTestDrive} esercita il percorso principale su H2; \code{RentalWorkflowTest} verifica le regole dell'entità. Non è presente una prova funzionale specifica per ciascuna alternativa o per richieste concorrenti.\\\hline
Flusso Principale &
\begin{enumerate}[label=\arabic*.,start=1,leftmargin=*,nosep]
  \item Il Customer seleziona un veicolo dal catalogo e apre la scheda dei test drive.
  \item Sceglie una data e uno degli orari proposti, dalle 09:00 alle 16:00.
  \item Il sistema verifica che il veicolo sia idoneo, che il Customer abbia meno di tre test drive futuri e che il veicolo non sia noleggiato nella data scelta.
  \item Il sistema verifica che lo slot di un'ora non sia già occupato per quel veicolo.
  \item Il sistema salva la prenotazione e ne comunica l'esito al Customer.
\end{enumerate}\\\hline
Flussi Alternativi &
\textbf{A1 (passo 2). Data o orario non validi:} se lo slot è nel passato o fuori dalla fascia consentita, il sistema rifiuta la prenotazione.\newline\newline
\textbf{A2 (passo 3). Veicolo non idoneo o limite raggiunto:} un veicolo non in vendita né a noleggio, oppure tre test drive futuri già registrati, impedisce il salvataggio.\newline\newline
\textbf{A3 (passi 3--4). Veicolo indisponibile:} un noleggio nel giorno scelto o un altro test drive nello stesso slot produce un conflitto.\newline\newline
\textbf{A4 (passo 4). Richieste concorrenti:} il controllo dello slot precede l'inserimento; il database non impone un vincolo univoco sullo slot, quindi due richieste simultanee potrebbero entrambe riuscire. Inoltre il servizio riconduce i minuti forniti all'inizio dell'ora: un input diretto come 10:30 è trattato come 10:00.\\\hline
\end{xltabular}
